\uuid{n8TJ}
\exo7id{4090}
\titre{exo7 4090}
\auteur{quercia}
\organisation{exo7}
\datecreate{2010-03-11}
\isIndication{false}
\isCorrection{true}
\chapitre{Equation différentielle}
\sousChapitre{Equations différentielles linéaires}
\cpge{mpsi-mp2i.derivabilite.d, pcsi.derivabilite.d, ptsi.derivabilite.d, mpsi-mp2i.analyse-asymptotique.a, pcsi.analyse-asymptotique.a, ptsi.analyse-asymptotique.a, bcpst1.equations-differentielles-lineaires.b, bcpst1.limites-continuite.a, bcpst2.integrales-generalisees}
\module{Analyse}
\niveau{L2}
\difficulte{}

\contenu{
\question{
Soit $f : \R \to\R$ de classe $\mathcal{C}^2$ telle que
$f''(t) + f'(t) + f(t) \to 0$ lorsque $t\to+\infty$. Démontrer que $f(t) \to 0$ lorsque $t\to\infty$.
}
\reponse{
On pose $\varphi(t) = f''(t)+f'(t)+f(t)$.

$f(t) = e^{-t/2} \left[ -\frac 2{\sqrt3}
         \int_{u=0}^t \varphi(u)e^{u/2} \sin\left({\frac {\sqrt3\,(u-t)}2 }\right), d u
        + A\cos\left(\frac {\sqrt3\,u}2 \right) + B\sin\left(\frac {\sqrt3\,u}2\right)
        \right ]$.
}
}
